\documentclass{ip-lab}

\lablevel{n3}
\labnumber{Laboratorio 2}
\labtitle{Ciclos sobre cadenas de caracteres}

\makeheader

\begin{document}

\maketitle

\begin{sectionbox}{Preparación del ambiente de trabajo}
  Cree un nuevo módulo de funciones en el que escribirá las funciones correspondientes a las actividades 1 a 3 de este laboratorio.
\end{sectionbox}


\begin{sectionbox}{Actividad 1: Última aparición de cada carácter}
Escriba una función que reciba como parámetro una cadena de caracteres y retorne una nueva cadena que contenga únicamente la última aparición de cada carácter, preservando el orden en que aparecen en la cadena original.

E.g., si la cadena de entrada es \texttt{\textquotedbl banana\textquotedbl}, la función debe retornar \texttt{\textquotedbl bna\textquotedbl}; si es \texttt{\textquotedbl trescientos\textquotedbl}, debe retornar \texttt{\textquotedbl rcientos\textquotedbl}.


\end{sectionbox}

\begin{sectionbox}{Actividad 2: Palabra más frecuente}
Escriba una función que reciba como parámetro una cadena de caracteres y retorne la palabra que aparece con mayor frecuencia en la cadena.

La función debe ignorar diferencias entre mayúsculas y minúsculas (i.e., \texttt{\textquotedbl HolA\textquotedbl} y \texttt{\textquotedbl hola\textquotedbl} se consideran la misma palabra). Si hay empate, puede retornar cualquiera de las palabras más frecuentes.

E.g., si la cadena de entrada es \texttt{\textquotedbl hola mundo Hola adiós HolA mundo\textquotedbl}, un ejemplo de retorno correcto es \texttt{\textquotesingle hola\textquotesingle}.
\end{sectionbox}

\pagebreak

\begin{sectionbox}{Actividad 3: Simetría de vocales}
Escriba una función que reciba como parámetro una cadena de caracteres y determine si las vocales en la cadena se leen igual de izquierda a derecha que de derecha a izquierda. No distinga entre mayúsculas y minúsculas. Presuma que no hay tildes en la cadena.

E.g., la función debe retornar \texttt{True} para la cadena de entrada \texttt{\textquotedbl Camina hacia la cima\textquotedbl}, pues las vocales se leen como \texttt{a, i, a, a, i, a, a, i, a} en ambas direcciones
\end{sectionbox}

\begin{sectionbox}{Entrega}
Entregue el archivo a través de Bloque Neón en el laboratorio del Nivel 3 designado como ``N3-L2: Ciclos sobre cadenas de caracteres''.
\end{sectionbox}

\end{document}
